% Generated draft: requires mathematical and source review.
\documentclass[11pt]{article}
\usepackage[margin=25mm]{geometry}
\usepackage{fontspec}
\setmainfont{Latin Modern Roman}
\usepackage{amsmath,amssymb,mathtools,unicode-math}
\setmathfont{Latin Modern Math}
\usepackage{enumitem,needspace,xurl,hyperref}
\setlength{\parindent}{0pt}
\setlength{\parskip}{5pt}
\setlength{\emergencystretch}{3em}
\newcommand{\N}{\mathbb{N}}
\newcommand{\Z}{\mathbb{Z}}
\newcommand{\Q}{\mathbb{Q}}
\newcommand{\R}{\mathbb{R}}
\newcommand{\C}{\mathbb{C}}
\newcommand{\F}{\mathbb{F}}
\newcommand{\A}{\mathbb{A}}
\newcommand{\PP}{\mathbb{P}}
\newcommand{\E}{\mathbb{E}}
\newcommand{\eps}{\varepsilon}
\newcommand{\dd}{\,\mathrm{d}}
\begin{document}
\section*{UnsolvedMath ID: 30001466}
\textbf{Source problem code:} catalogue entry
\subsection{Definitions and mathematical statement}
Use the standard convention that a \emph{limit group} is a finitely generated fully residually free group: a finitely generated group \(G\) such that for every finite subset \(S\subseteq G\setminus\{1_G\}\), there is a homomorphism \(\phi:G\to F_2\) satisfying \(\phi(s)\ne 1_{F_2}\) for every \(s\in S\), where \(F_2\) is the free group on two generators.

Two countable groups \(G\) and \(H\) are \emph{measure equivalent} if there exists a nonzero standard \(\sigma\)-finite measure space \((\Omega,\mu)\) with commuting, measure-preserving, essentially free actions of \(G\) and \(H\), each admitting a measurable fundamental domain of finite positive measure. Here essential freeness means that almost every point has trivial stabilizer for the respective action. A fundamental domain for the \(G\)-action is a measurable set \(D_G\subseteq\Omega\) whose translates \(\{gD_G:g\in G\}\) partition \(\Omega\) up to null sets; the definition for \(H\) is analogous.

For an at most countable set \(B\), write \(F(B)\) for the free group with basis \(B\), represented by reduced words in the symbols \(B\cup B^{-1}\). The source does not restrict the rank of the target free group. Accordingly, the formulation below permits any at most countable basis, including the empty basis and a one-element basis; any intended restriction to nonabelian or finite-rank targets would require clarification.

\textbf{Question.} Is it true that for every limit group \(G\), there exists an at most countable set \(B\) such that \(G\) and \(F(B)\) are measure equivalent?
\[
\forall G\;\bigl(G\text{ is a limit group}\ \Longrightarrow\ \exists B\text{ at most countable}:\ G\text{ is measure equivalent to }F(B)\bigr).
\]
The supplied record does not explicitly specify conventions for trivial or abelian limit groups. The fully residually free convention stated above includes them; any intended exclusion remains to be checked against the original report.
\subsection{Short English statement}
Does every finitely generated fully residually free group admit a measure-equivalence coupling with some free group? Such a coupling consists of commuting essentially free measure-preserving actions on one measure space, with a finite positive measure fundamental domain for each action. The supplied statement leaves the target free group's rank unspecified.
\subsection{Sources}
\begin{itemize}
\item Local UnsolvedMath registry, permanent ID \(30001466\), source problem code \texttt{catalogue entry}, title \emph{Measure Equivalence of Limit and Free Groups}. The supplied record attributes the question to the 2010 Oberwolfach Reports collection and names no individual proposer. It labels the extracted statement exact and reports a prior manual statement check; these are supplied registry qualifications.
\item \emph{Non-positive Curvature and Geometric Structures in Group Theory} (2010), Oberwolfach Reports, pp.\ 1165--1224, DOI \texttt{10.4171/owr/2010/20}. Supplied original-report reference: \url{https://doi.org/10.4171/owr/2010/20}.
\item Bridson--Tweedale--Wilton, supplied literature reference: \url{https://arxiv.org/abs/math/0511275}. The record supplies no title or specific theorem locator for this reference.
\item Supplied literature assessment, checked \texttt{2026-08-21}: the general question is labeled open; the record reports that primary work establishes the result for elementarily free groups and provides only partial results for arbitrary limit groups. This is supplied dated curation derived from the Oberwolfach material, not an independently verified literature assessment. Its status and cited claims require independent verification.
\end{itemize}
\end{document}
